% Original combinatorial signs bridge.
\subsection*{Worked calculation: tracking fermionic signs in three orbitals}
The antisymmetric Fock algebra is sometimes difficult because
operators act \emph{from right to left} while the canonical
occupation convention orders creation operators from left to right.
We calculate a three-mode example explicitly.

Choose orbitals $1<2<3$, vacuum $\ket0$, and
\begin{equation}
\ket{1,0,1}=\op a_1^\dagger\op a_3^\dagger\ket0.
\end{equation}
Act with $\op a_2^\dagger$. It must cross the occupied
lower-index mode $1$ to restore canonical order:
\begin{align}
\op a_2^\dagger\ket{1,0,1}
&=\op a_2^\dagger\op a_1^\dagger\op a_3^\dagger\ket0\notag\\
&=-\op a_1^\dagger\op a_2^\dagger\op a_3^\dagger\ket0
=\boxed{-\ket{1,1,1}}.
\end{align}
The minus sign is exactly $(-1)^{\nu_2}$, since one
lower-index orbital is occupied. Acting instead with
$\op a_1$ on $\ket{1,0,1}$,
\begin{equation}
\op a_1\op a_1^\dagger\op a_3^\dagger\ket0
=(1-\op a_1^\dagger\op a_1)\op a_3^\dagger\ket0
=\ket{0,0,1},
\end{equation}
because $\op a_1\ket0=0$.
This result has positive sign: mode 1 has no occupied
lower-index orbitals.

\paragraph{A complete one-body transition.}
Evaluate $\op a_2^\dagger\op a_3$ on
$\ket{1,0,1}$ by applying the rightmost operator first.
Since $\op a_3$ must pass the particle in mode 1,
\begin{equation}
\op a_3\ket{1,0,1}
=\op a_3\op a_1^\dagger\op a_3^\dagger\ket0
=-\op a_1^\dagger(1-\op a_3^\dagger\op a_3)\ket0
=-\ket{1,0,0}.
\end{equation}
Now create in mode 2:
\begin{equation}
\op a_2^\dagger\op a_3\ket{1,0,1}
=-\op a_2^\dagger\op a_1^\dagger\ket0
=+\op a_1^\dagger\op a_2^\dagger\ket0
=\boxed{+\ket{1,1,0}}.
\end{equation}
Two minus signs cancel. This small calculation is the
concrete meaning of the matrix element
$t_{23}\op a_2^\dagger\op a_3$
in a one-body Hamiltonian.

\paragraph{Where does the two-body factor of one half come from?}
Start with the first-quantized sum over unordered pairs
$\op V=\sum_{\alpha<\beta}v(\vect r_\alpha,\vect r_\beta)$.
For a symmetric interaction
$v(\vect r_\alpha,\vect r_\beta)
=v(\vect r_\beta,\vect r_\alpha)$,
\begin{equation}
\sum_{\alpha<\beta}v_{\alpha\beta}
=\frac12\sum_{\alpha\neq\beta}v_{\alpha\beta}.
\end{equation}
The right-hand sum counts each pair twice. After expanding
the one-particle coordinate basis and using creation
and annihilation operators, this becomes
\begin{equation}
\op V=\frac12\sum_{ijkl}
V_{ij;kl}\op a_i^\dagger\op a_j^\dagger
\op a_l\op a_k.
\end{equation}
The prefactor follows from \emph{pair counting}; anticommutation
then supplies exchange signs. If antisymmetrized two-body
matrix elements are used instead of the unsymmetrized
$V_{ij;kl}$ defined above, the numerical prefactor must
be reconsidered to avoid double counting.

\paragraph{Self-check.}
Calculate $\op a_3^\dagger\op a_2\ket{1,1,0}$
and confirm that the outcome is $+\ket{1,0,1}$,
the Hermitian-adjoint transition of the preceding example.
